%===============================================================================
% БАКАЛАВРЫН ТӨГСӨЛТИЙН АЖЛЫН ИЛТГЭЛ — М.Түвшинжаргал В221940619
% ШУТИС-МХТС, Мэдээллийн технологийн тэнхим, 2026
%===============================================================================
\documentclass[14pt,aspectratio=169]{beamer}

%--- Encoding & fonts ----------------------------------------------------------
\usepackage[T2A]{fontenc}
\usepackage[utf8]{inputenc}
\usepackage[mongolian,english]{babel}

%--- Packages ------------------------------------------------------------------
\usepackage{tikz}
\usetikzlibrary{arrows.meta,shapes.geometric,calc}
\usepackage{booktabs}
\usepackage{xcolor}
\usepackage{amsmath,amssymb}
\usepackage{array,colortbl}
\usepackage{multirow}

%--- Custom colors -------------------------------------------------------------
\definecolor{navyblue}{RGB}{8,40,85}
\definecolor{navylight}{RGB}{18,60,120}
\definecolor{goldcolor}{RGB}{210,155,20}
\definecolor{lightbg}{RGB}{235,242,255}
\definecolor{okgreen}{RGB}{20,140,60}
\definecolor{badred}{RGB}{190,30,30}
\definecolor{partcolor}{RGB}{200,130,0}
\definecolor{rowgray}{RGB}{245,247,250}

%--- Beamer theme --------------------------------------------------------------
\usetheme{default}
\usecolortheme{default}

\setbeamercolor{frametitle}{bg=navyblue,fg=white}
\setbeamercolor{title}{fg=navyblue}
\setbeamercolor{subtitle}{fg=navylight}
\setbeamercolor{author}{fg=navyblue}
\setbeamercolor{institute}{fg=navyblue!80}
\setbeamercolor{date}{fg=navyblue!70}
\setbeamercolor{block title}{bg=navyblue,fg=white}
\setbeamercolor{block body}{bg=lightbg,fg=black}
\setbeamercolor{alerted text}{fg=goldcolor!85!black}
\setbeamercolor{structure}{fg=navyblue}
\setbeamercolor{itemize item}{fg=navyblue}
\setbeamercolor{itemize subitem}{fg=navylight}
\setbeamercolor{enumerate item}{fg=navyblue}
\setbeamercolor{section in toc}{fg=navyblue}

%--- Frame title with gold underline -------------------------------------------
\setbeamertemplate{frametitle}{%
  \nointerlineskip
  \begin{beamercolorbox}[wd=\paperwidth,ht=2.6ex,dp=0.8ex,leftskip=1.2em,
                         rightskip=1.2em]{frametitle}
    \usebeamerfont{frametitle}\insertframetitle
  \end{beamercolorbox}%
  \nointerlineskip
  {\color{goldcolor}\hrule height 1.6pt}%
}

%--- Navigation off, custom footer --------------------------------------------
\setbeamertemplate{navigation symbols}{}
\setbeamertemplate{footline}{%
  \leavevmode%
  \hbox{%
    \begin{beamercolorbox}[wd=.38\paperwidth,ht=2.2ex,dp=1ex,
                           leftskip=0.6em]{block title}%
      \tiny М.Түвшинжаргал \textbullet\ В221940619
    \end{beamercolorbox}%
    \begin{beamercolorbox}[wd=.44\paperwidth,ht=2.2ex,dp=1ex,center]{block body}%
      \tiny Спорт заалны цаг захиалгын 3 давхаргат систем
    \end{beamercolorbox}%
    \begin{beamercolorbox}[wd=.18\paperwidth,ht=2.2ex,dp=1ex,
                           rightskip=0.6em,right]{block title}%
      \tiny \insertframenumber\,/\,\inserttotalframenumber
    \end{beamercolorbox}%
  }%
}

%--- Bullet styles -------------------------------------------------------------
\setbeamertemplate{itemize item}{\color{navyblue}\rule[0.3ex]{0.6ex}{0.6ex}}
\setbeamertemplate{itemize subitem}{\color{goldcolor}$\triangleright$}
\setbeamertemplate{enumerate item}{\color{navyblue}\textbf{\insertenumlabel.}}

%--- Custom title page ---------------------------------------------------------
\setbeamertemplate{title page}{%
  \vspace{-0.4cm}
  {\color{navyblue}\hrule height 2.5pt}
  \vspace{0.25cm}
  \begin{center}
    {\footnotesize\color{navyblue!70}
     ШУТИС — Мэдээлэл Харилцаа Технологийн Сургууль\\
     Мэдээллийн технологийн тэнхим \quad|\quad 2026 он}
  \end{center}
  {\color{goldcolor}\hrule height 1pt}
  \vspace{0.3cm}
  \begin{center}
    {\large\bfseries\color{navyblue}БАКАЛАВРЫН ТӨГСӨЛТИЙН АЖЛЫН ИЛТГЭЛ}
  \end{center}
  \vspace{0.2cm}
  \begin{beamercolorbox}[sep=0.4em,center,rounded=true,shadow=false]{block title}
    \usebeamerfont{title}\large\inserttitle
  \end{beamercolorbox}
  \vspace{0.35cm}
  \begin{columns}[T]
    \begin{column}{0.5\textwidth}
      \centering\small
      \textbf{Оюутан:} М.Түвшинжаргал\\
      \textcolor{navyblue!70}{В221940619}
    \end{column}
    \begin{column}{0.5\textwidth}
      \centering\small
      \textbf{Удирдагч:} ахлах багш, доктор Д.Эрдэнэтуяа\\
      \textbf{Зөвлөх:} багш, доктор Б.Намжилдорж
    \end{column}
  \end{columns}
  \vspace{0.3cm}
  {\color{navyblue}\hrule height 2.5pt}
}

%--- Title info ----------------------------------------------------------------
\title[Спорт заалны цаг захиалгын систем]{%
  Үүлэн орчинд суурилсан спорт заалны цаг захиалгын\\
  3 давхаргат системийн загварчлал, хэрэгжилт}
\author{}
\institute{}
\date{}

%--- Helper commands -----------------------------------------------------------
\newcommand{\YES}{\textcolor{okgreen}{\checkmark}}
\newcommand{\NO}{\textcolor{badred}{$\times$}}
\newcommand{\PART}{\textcolor{partcolor}{$\triangle$}}

%===============================================================================
\begin{document}
%===============================================================================

%--- 01 Title ------------------------------------------------------------------
\begin{frame}[plain]
  \titlepage
\end{frame}

%--- 02 Outline ----------------------------------------------------------------
\begin{frame}{Илтгэлийн агуулга}
  \begin{columns}[T]
    \begin{column}{0.48\textwidth}
      \tableofcontents[sections={1-5}]
    \end{column}
    \begin{column}{0.48\textwidth}
      \tableofcontents[sections={6-9}]
    \end{column}
  \end{columns}
\end{frame}

%===============================================================================
\section{Асуудлын тодорхойлолт}
%===============================================================================
\begin{frame}{Асуудлын тодорхойлолт}
  \begin{columns}[T]
    \begin{column}{0.46\textwidth}
      \begin{block}{Одоогийн нөхцөл байдал}
        \begin{itemize}
          \item Даланзадгад хотод \textbf{5 спорт заал}
          \item Иргэд зөвхөн \textbf{утсаар} эсвэл \textbf{биечлэн} цаг захиалдаг
          \item Шугам завгүй, буруу ойлголцол нийтлэг
          \item Боломжтой цагийн мэдээлэл \textbf{ил тод бус}
          \item Захиалгын баталгаажуулалт \textbf{байхгүй}
        \end{itemize}
      \end{block}
    \end{column}
    \begin{column}{0.50\textwidth}
      \begin{alertblock}{Шийдвэрлэх асуудлууд}
        \begin{enumerate}
          \item Цаг захиалгыг \textbf{цахимжуулах}
          \item Бодит цагийн хуваарь харуулах
          \item Захиалгын автомат баталгаажуулалт
          \item Байгууллагын \textbf{тогтмол} захиалгыг дэмжих
          \item Хэрэглэгчдэд \textbf{AI туслагч} санал болгох
        \end{enumerate}
      \end{alertblock}
      \vspace{4pt}
      \begin{beamercolorbox}[sep=4pt,rounded=true]{block body}
        \centering\small
        \textbf{Зорилтот хэрэглэгчид:} Даланзадгад хотын иргэд
        болон спорт заалны ажилтнууд
      \end{beamercolorbox}
    \end{column}
  \end{columns}
\end{frame}

%===============================================================================
\section{Судлагдсан байдал}
%===============================================================================
\begin{frame}{Ижил төстэй системүүдийн харьцуулалт}
  \begin{center}
  \small
  \renewcommand{\arraystretch}{1.25}
  \begin{tabular}{lcccc}
    \toprule
    \rowcolor{navyblue}
    {\color{white}\textbf{Онцлог}} &
    {\color{white}\textbf{Mindbody}} &
    {\color{white}\textbf{Vagaro}} &
    {\color{white}\textbf{Нутгийн апп}} &
    {\color{white}\textbf{Манай систем}} \\
    \midrule
    \rowcolor{rowgray}
    Монгол хэл / OTP нэвтрэлт & \NO & \NO & \PART & \YES \\
    Бодит цагийн хуваарь       & \YES & \YES & \PART & \YES \\
    \rowcolor{rowgray}
    Байгууллагын тогтмол захиалга & \PART & \PART & \NO & \YES \\
    AI чатбот туслагч          & \NO & \NO & \NO & \YES \\
    \rowcolor{rowgray}
    Cloud auto-deploy (CI/CD)  & N/A & N/A & \NO & \YES \\
    Орон нутгийн хэрэгцээ      & \NO & \NO & \PART & \YES \\
    \rowcolor{rowgray}
    Үнэ (сард)                 & \$129+ & \$25+ & Үнэгүй & Үнэгүй \\
    \bottomrule
  \end{tabular}
  \end{center}
  \vspace{4pt}
  \begin{beamercolorbox}[sep=4pt,center,rounded=true]{block body}
    \small Олон улсын платформууд нь Монголын орон нутгийн хэрэгцээнд
    \textbf{тохиромжгүй}. Манай систем нь тусгай шийдлийг санал болгодог.
  \end{beamercolorbox}
\end{frame}

\begin{frame}{Технологийн сонголт — Flutter}
  \begin{center}
  \small
  \renewcommand{\arraystretch}{1.3}
  \begin{tabular}{lccc}
    \toprule
    \rowcolor{navyblue}
    {\color{white}\textbf{Шалгуур}} &
    {\color{white}\textbf{Flutter}} &
    {\color{white}\textbf{React Native}} &
    {\color{white}\textbf{Native (Java/Swift)}} \\
    \midrule
    \rowcolor{rowgray}
    Android + iOS нэг кодоос & \YES & \YES & \NO (2 апп) \\
    Гүйцэтгэл (performance) & \textbf{Маш сайн} & Сайн & Маш сайн \\
    \rowcolor{rowgray}
    Firebase SDK дэмжлэг & \YES & \YES & \YES \\
    Хөгжүүлэлтийн хугацаа & \textbf{1×} & 1× & 2× \\
    \rowcolor{rowgray}
    Hot reload & \YES & \YES & \NO \\
    Material Design 3 & \textbf{Нативаар} & Хязгаарлагдмал & Хязгаарлагдмал \\
    \rowcolor{rowgray}
    Монгол хэлний дэмжлэг & \YES & \YES & \YES \\
    \bottomrule
  \end{tabular}
  \end{center}
  \vspace{4pt}
  \begin{beamercolorbox}[sep=4pt,rounded=true]{block body}
    \small\centering
    \textbf{Flutter} сонгосон шалтгаан: Нэг кодоос хоёр платформд гаргах,
    гүйцэтгэл өндөр, Firebase-тэй нативаар нэгдсэн
  \end{beamercolorbox}
\end{frame}

%===============================================================================
\section{Системийн архитектур}
%===============================================================================
\begin{frame}{3 Давхаргат системийн архитектур}
  \begin{center}
  \begin{tikzpicture}[scale=0.88]
    % Styles
    \tikzset{
      cbox/.style={rectangle, rounded corners=5pt, draw, text width=2.7cm,
                   align=center, minimum height=1.25cm, font=\small},
      myarr/.style={-{Stealth[length=5pt,width=4pt]}, thick}
    }
    % Layer labels on left
    \node[font=\footnotesize, rotate=90, navyblue, text width=1.5cm, align=center]
      at (-1.8, 0.6) {\textbf{1-р давхар}\\Клиент};
    \node[font=\footnotesize, rotate=90, navyblue, text width=1.5cm, align=center]
      at (-1.8, -2.0) {\textbf{2,3-р давхар}\\Сервер};

    % Client layer
    \node[cbox, fill=yellow!20, draw=yellow!60!black] (app) at (0.2, 0.6)
      {\textbf{Flutter Апп}\\Android / iOS};

    % Backend layer
    \node[cbox, fill=orange!20, draw=orange!60] (ec2) at (4.6, 0.6)
      {\textbf{Node.js/Express}\\AWS EC2 (Tokyo)};
    \node[cbox, fill=blue!15, draw=blue!50] (fb) at (0.2, -2.0)
      {\textbf{Firebase}\\Auth + Firestore\\+ Cloud Functions};
    \node[cbox, fill=green!15, draw=green!50] (ai) at (8.4, 0.6)
      {\textbf{AI үйлчилгээ}\\Claude Haiku\\+ Groq Llama};
    \node[cbox, fill=gray!15, draw=gray!60, minimum height=0.9cm] (ecr)
      at (4.6, -2.0) {\textbf{AWS ECR}\\Docker Registry};
    \node[cbox, fill=gray!15, draw=gray!60, minimum height=0.9cm] (gh)
      at (2.0, -2.0) {\textbf{GitHub Actions}\\CI/CD};

    % Arrows
    \draw[myarr, navyblue] (app.east) -- node[above,font=\tiny]{REST API} (ec2.west);
    \draw[myarr, blue!60] (app.south) -- node[left,font=\tiny]{SDK} (fb.north);
    \draw[myarr, green!60] (ec2.east) -- node[above,font=\tiny]{AI API} (ai.west);
    \draw[myarr, gray!70] (gh.east) -- node[below,font=\tiny]{push image} (ecr.west);
    \draw[myarr, gray!70] (ecr.north) -- node[right,font=\tiny]{pull} (ec2.south);

    % Dashed box for infra
    \draw[dashed, navyblue!40, rounded corners=6pt]
      (-0.3,-1.3) rectangle (9.2, 1.5);
    \node[font=\tiny, navyblue!50] at (8.5, 1.35) {Системийн хил};
  \end{tikzpicture}
  \end{center}
\end{frame}

%===============================================================================
\section{UML диаграммууд}
%===============================================================================
\begin{frame}{Use Case диаграмм}
  \begin{center}
  \begin{tikzpicture}[scale=0.82]
    \tikzset{
      uc/.style={ellipse, draw=navyblue!70, fill=blue!8,
                 text width=2.6cm, align=center, font=\tiny,
                 minimum height=0.65cm},
      ucarr/.style={navyblue!60, thin}
    }
    % System boundary
    \draw[navyblue!50, thick, rounded corners=6pt]
      (1.5,-0.6) rectangle (10.5,5.4);
    \node[font=\footnotesize, navyblue] at (6.0,5.15) {Спорт заалны цаг захиалгын систем};

    % Actor — stick figure
    \draw[thick] (-0.4, 3.8) circle (0.28);
    \draw[thick] (-0.4, 3.5) -- (-0.4, 2.6);
    \draw[thick] (-0.9, 3.2) -- (0.1, 3.2);
    \draw[thick] (-0.4, 2.6) -- (-0.8, 2.0);
    \draw[thick] (-0.4, 2.6) -- (0.0, 2.0);
    \node[below, font=\scriptsize] at (-0.4, 1.95) {\textbf{Хэрэглэгч}};

    % Use cases (2 columns)
    \node[uc] (uc1) at (3.5, 4.8) {OTP-р нэвтрэх};
    \node[uc] (uc2) at (3.5, 3.7) {Хуваарь харах};
    \node[uc] (uc3) at (3.5, 2.6) {Цаг захиалах};
    \node[uc] (uc4) at (3.5, 1.5) {Захиалга цуцлах};
    \node[uc] (uc5) at (3.5, 0.4) {Профайл харах};
    \node[uc] (uc6) at (7.5, 4.8) {AI чат ашиглах};
    \node[uc] (uc7) at (7.5, 3.5) {Мэдэгдэл хүлээн авах};
    \node[uc] (uc8) at (7.5, 2.2) {Захиалгын код харах};
    \node[uc] (uc9) at (7.5, 0.9) {Захиалгын түүх};

    % Lines from actor
    \foreach \uc in {uc1,uc2,uc3,uc4,uc5} {
      \draw[ucarr] (0.1, 2.9) -- (\uc.west);
    }
    \foreach \uc in {uc6,uc7,uc8,uc9} {
      \draw[ucarr] (0.1, 2.9) -- (\uc.west);
    }
  \end{tikzpicture}
  \end{center}
\end{frame}

\begin{frame}{CI/CD хоолойн диаграмм}
  \begin{center}
  \begin{tikzpicture}[scale=0.85]
    \tikzset{
      stepa/.style={-{Stealth[length=5pt]}, thick, gray!70},
      stepbox/.style={rectangle, rounded corners=4pt, text width=1.9cm,
                      align=center, minimum height=1.1cm, font=\small}
    }
    \node[stepbox, draw=blue!70, fill=blue!12] (dev)
      at (0.0, 0) {\textbf{Хөгжүүлэгч}\\git push};
    \node[stepbox, draw=teal!70, fill=teal!12] (gh)
      at (2.8, 0) {\textbf{GitHub}\\Trigger};
    \node[stepbox, draw=orange!70, fill=orange!12] (build)
      at (5.6, 0) {\textbf{Docker}\\Build};
    \node[stepbox, draw=red!70, fill=red!10] (ecr)
      at (8.4, 0) {\textbf{AWS ECR}\\Push image};
    \node[stepbox, draw=navyblue!70, fill=navyblue!12] (ec2)
      at (11.2,0) {\textbf{AWS EC2}\\Deploy};

    \draw[stepa] (dev)   -- node[above,font=\tiny]{backend/} (gh);
    \draw[stepa] (gh)    -- node[above,font=\tiny]{Actions} (build);
    \draw[stepa] (build) -- node[above,font=\tiny]{SHA tag} (ecr);
    \draw[stepa] (ecr)   -- node[above,font=\tiny]{SSH} (ec2);

    % Time label
    \draw[<->, gray!60, dashed] (0,-0.9) -- node[below,font=\tiny,gray]{3--4 минут} (11.2,-0.9);

    % Step numbers
    \node[font=\tiny, gray!80] at (0,   -0.5) {(1)};
    \node[font=\tiny, gray!80] at (2.8, -0.5) {(2)};
    \node[font=\tiny, gray!80] at (5.6, -0.5) {(3)};
    \node[font=\tiny, gray!80] at (8.4, -0.5) {(4)};
    \node[font=\tiny, gray!80] at (11.2,-0.5) {(5)};
  \end{tikzpicture}
  \end{center}
  \vspace{4pt}
  \begin{beamercolorbox}[sep=4pt,center,rounded=true]{block body}
    \small \texttt{git push} $\Rightarrow$ GitHub Actions автоматаар ажиллаж
    $\Rightarrow$ EC2 дээрх контейнер \textbf{хүний гараас хамааралгүйгээр} шинэчлэгдэнэ
  \end{beamercolorbox}
\end{frame}

%===============================================================================
\section{Өгөгдлийн сангийн загвар}
%===============================================================================
\begin{frame}{Firestore өгөгдлийн сангийн схем}
  \begin{center}
  \begin{tikzpicture}[scale=0.88]
    \tikzset{
      coll/.style={rectangle, draw=navyblue!70, fill=blue!8,
                   rounded corners=4pt, minimum width=3.4cm},
      field/.style={font=\tiny, navyblue!80},
      pk/.style={font=\tiny\bfseries, navyblue},
      fk/.style={font=\tiny, blue!60},
      rel/.style={-{Stealth[length=4pt]}, thick, navyblue!50, dashed}
    }
    % users collection
    \node[coll, minimum height=2.6cm] (users) at (0, 0) {};
    \node[font=\small\bfseries, navyblue] at (0, 1.0) {\texttt{users}};
    \draw[navyblue!40] (-1.7, 0.65) -- (1.7, 0.65);
    \node[pk] at (0, 0.4) {uid (doc ID, PK)};
    \node[field] at (0, 0.05) {name : string};
    \node[field] at (0,-0.3)  {phone : string};
    \node[field] at (0,-0.65) {createdAt : timestamp};

    % bookings collection
    \node[coll, minimum height=4.2cm] (bk) at (5.5, -0.3) {};
    \node[font=\small\bfseries, navyblue] at (5.5, 1.65) {\texttt{bookings}};
    \draw[navyblue!40] (3.8, 1.3) -- (7.2, 1.3);
    \node[pk]    at (5.5, 0.95) {bookingCode (PK)};
    \node[fk]    at (5.5, 0.6)  {userId (FK $\to$ users)};
    \node[fk]    at (5.5, 0.25) {venueId (FK $\to$ venues)};
    \node[field] at (5.5,-0.1)  {venueName : string};
    \node[field] at (5.5,-0.45) {date : string};
    \node[field] at (5.5,-0.8)  {timeSlot : string};
    \node[field] at (5.5,-1.15) {status : active|completed|cancelled};
    \node[field] at (5.5,-1.5)  {createdAt : timestamp};

    % fixed_bookings collection
    \node[coll, minimum height=3.2cm] (fb) at (11.0, -0.3) {};
    \node[font=\small\bfseries, navyblue] at (11.0, 1.2) {\texttt{fixed\_bookings}};
    \draw[navyblue!40] (9.3, 0.85) -- (12.7, 0.85);
    \node[fk]    at (11.0, 0.5)  {venueId (FK $\to$ venues)};
    \node[field] at (11.0, 0.15) {organizationName : string};
    \node[field] at (11.0,-0.2)  {dayOfWeek : 0--6};
    \node[field] at (11.0,-0.55) {timeSlot : string};
    \node[field] at (11.0,-0.9)  {isActive : boolean};

    % Relationships
    \draw[rel] (1.7, 0.2) -- node[above,font=\tiny]{1:N} (3.8, 0.7);
    \draw[rel] (7.2, 0.0) -- node[above,font=\tiny]{venueId} (9.3, 0.2);

    % venues note
    \node[draw=gray!50, fill=gray!8, rounded corners=3pt, font=\tiny,
          text width=2.4cm, align=center] (vn) at (5.5, -2.6)
      {\textbf{venues}\\\small Flutter кодод\\hardcoded (5 заал)};
    \draw[gray!50, dashed, thin] (bk.south) -- (vn.north);
    \draw[gray!50, dashed, thin] (fb.south) -- (vn.east);
  \end{tikzpicture}
  \end{center}
\end{frame}

%===============================================================================
\section{Хэрэгжүүлэлт}
%===============================================================================
\begin{frame}{Flutter апп — UI/UX ба гол функцүүд}
  \begin{columns}[T]
    \begin{column}{0.30\textwidth}
      % Phone mockup (TikZ)
      \centering
      \begin{tikzpicture}[scale=0.78]
        % Phone body
        \draw[rounded corners=8pt, fill=black!88, draw=black!70, line width=1.5pt]
          (0,0) rectangle (3.0, 5.6);
        % Screen area
        \fill[black!5] (0.18,0.4) rectangle (2.82, 5.15);
        % Notch
        \fill[black!88] (1.0,5.0) rectangle (2.0, 5.2);
        % Status bar color
        \fill[black!80] (0.18,4.8) rectangle (2.82,5.15);
        % App title bar
        \fill[black!75] (0.18,4.35) rectangle (2.82,4.8);
        \node[font=\tiny, white, align=center] at (1.5, 4.57)
          {Dalandzadgad Sport};
        % Background color (desert dark)
        \fill[black!70] (0.18,0.4) rectangle (2.82,4.35);
        % Cards
        \fill[black!55, rounded corners=3pt] (0.28,3.65) rectangle (2.72,4.28);
        \node[font=\tiny, yellow!80, align=center] at (1.5, 3.96)
          {Goviin Arena \hfill $\triangleright$};
        \fill[black!55, rounded corners=3pt] (0.28,2.90) rectangle (2.72,3.53);
        \node[font=\tiny, yellow!80, align=center] at (1.5, 3.21)
          {Omnogovi Sport Hall \hfill $\triangleright$};
        \fill[black!55, rounded corners=3pt] (0.28,2.15) rectangle (2.72,2.78);
        \node[font=\tiny, yellow!80] at (1.5, 2.46)
          {Goviin Volleyball};
        \fill[black!55, rounded corners=3pt] (0.28,1.40) rectangle (2.72,2.03);
        \node[font=\tiny, yellow!80] at (1.5, 1.71)
          {Stadium Volleyball};
        % Bottom nav
        \fill[black!80] (0.18,0.4) rectangle (2.82,1.1);
        \foreach \x/\lbl in {0.55/Home,1.15/Book,1.75/Chat,2.35/Me} {
          \node[font=\tiny, white] at (\x+0.2, 0.72) {\lbl};
        }
        % Home button
        \fill[gray!50] (1.5,0.2) circle (0.12);
      \end{tikzpicture}
    \end{column}
    \begin{column}{0.67\textwidth}
      \textbf{Dark Desert Theme}\\
      \small
      \begin{tabular}{ll}
        Цөлийн хар & \colorbox{black!85}{\color{white}\small\texttt{\#120D06}} \\
        Алтан элс  & \colorbox{yellow!70}{\small\texttt{\#F0B429}} \\
        Нарны жаргалт & \colorbox{orange!60}{\small\texttt{\#D4622A}}
      \end{tabular}
      \vspace{6pt}

      \textbf{Үндсэн функцүүд:}
      \begin{itemize}
        \item Firebase OTP утасны нэвтрэлт
        \item Бодит цагийн цагийн хуваарь (Firestore)
        \item \textbf{DBK-xxxxx} форматын захиалгын код
        \item fixed\_bookings — байгууллагын тогтмол цаг
        \item Захиалгаас \textbf{1 цагийн өмнө} локал сануулга
        \item Захиалга цуцлах / захиалгын түүх
        \item AI чатбот дэлгэц
      \end{itemize}
    \end{column}
  \end{columns}
\end{frame}

\begin{frame}{AI чатботын архитектур}
  \begin{center}
  \begin{tikzpicture}[scale=0.88]
    \tikzset{
      abox/.style={rectangle, rounded corners=5pt,
                   text width=3.2cm, align=center, minimum height=1.2cm, font=\small},
      aarr/.style={-{Stealth[length=5pt]}, thick}
    }
    % Flutter App
    \node[abox, draw=blue!70, fill=blue!12] (app) at (0,0)
      {\textbf{Flutter Апп}\\ChatScreen widget};

    % Two paths
    \node[abox, draw=teal!70, fill=teal!12, text width=3.5cm] (cf) at (5.0, 1.5)
      {\textbf{Firebase Cloud Functions}\\Claude Haiku 4.5\\(Anthropic API)};
    \node[abox, draw=orange!70, fill=orange!12, text width=3.5cm] (ex) at (5.0,-1.5)
      {\textbf{Node.js/Express}\\Groq Llama 3.1-8b\\(нөөц горим)};

    % Response merge
    \node[abox, draw=navyblue!70, fill=navyblue!12] (resp) at (10.2, 0)
      {\textbf{Хариулт}\\хэрэглэгчид\\харуулна};

    % Arrows
    \draw[aarr, teal!70] (app.east) |- node[above,font=\tiny,pos=0.8]{Үндсэн горим} (cf.west);
    \draw[aarr, orange!70] (app.east) |- node[below,font=\tiny,pos=0.8]{Нөөц горим} (ex.west);
    \draw[aarr, teal!70] (cf.east) -| (resp.north);
    \draw[aarr, orange!70] (ex.east) -| (resp.south);

    % Labels
    \node[font=\tiny, teal!80] at (5.0, 2.85)
      {System prompt: Спорт заалны туслагч};
    \node[font=\tiny, orange!80] at (5.0,-2.7)
      {Хэрэглэгч нэвтрэлтгүй эсвэл CF унтарсан үед};

    % Failover arrow
    \draw[dashed, gray!60, -{Stealth[length=4pt]}] (cf.south) --
      node[right,font=\tiny,gray]{failover} (ex.north);
  \end{tikzpicture}
  \end{center}
\end{frame}

%===============================================================================
\section{Туршилтын үр дүн}
%===============================================================================
\begin{frame}{Функциональ туршилтын үр дүн}
  \begin{columns}[T]
    \begin{column}{0.54\textwidth}
      \small
      \renewcommand{\arraystretch}{1.2}
      \begin{tabular}{clc}
        \toprule
        \rowcolor{navyblue}
        {\color{white}\textbf{№}} &
        {\color{white}\textbf{Тест тохиолдол}} &
        {\color{white}\textbf{Үр дүн}} \\
        \midrule
        \rowcolor{rowgray}
        1 & OTP нэвтрэлт, users бүртгэл & \YES \\
        2 & 5 спорт заал нүүр дэлгэцэд  & \YES \\
        \rowcolor{rowgray}
        3 & fixed\_bookings хуваарьт тусах & \YES \\
        4 & DBK-xxxxx захиалгын код        & \YES \\
        \rowcolor{rowgray}
        5 & 1 цагийн өмнө локал сануулга  & \YES \\
        6 & Захиалга цуцлах               & \YES \\
        \rowcolor{rowgray}
        7 & Cron — completed болгох        & \YES \\
        8 & Claude Haiku хариулт          & \YES \\
        \rowcolor{rowgray}
        9 & Groq Llama нөөц горим         & \YES \\
        10 & GitHub Actions auto-deploy   & \YES \\
        \bottomrule
        \multicolumn{3}{r}{\small\textbf{10 / 10 тэнцсэн}}
      \end{tabular}
    \end{column}
    \begin{column}{0.44\textwidth}
      \textbf{Гүйцэтгэлийн хэмжилт:}
      \vspace{6pt}

      \begin{tikzpicture}[scale=0.72]
        % Axes
        \draw[-{Stealth[length=4pt]}, thick] (-0.2,0) -- (6.2,0);
        \draw[-{Stealth[length=4pt]}, thick] (0,-0.1) -- (0,3.6)
          node[above,font=\tiny]{секунд};
        % Grid
        \foreach \y/\lbl in {1.0/1.0,2.0/2.0,3.0/3.0} {
          \draw[gray!25] (0,\y) -- (6.0,\y);
          \node[left,font=\tiny,gray!70] at (0,\y) {\lbl с};
        }
        % Target line
        \draw[dashed,red!50,thick] (0,3.0) -- (6.0,3.0)
          node[right,font=\tiny,red!50]{зорилт 3с};
        % Bars
        \fill[blue!60]    (0.2,0) rectangle (0.8,2.3);
        \fill[green!55]   (1.2,0) rectangle (1.8,0.6);
        \fill[orange!65]  (2.2,0) rectangle (2.8,0.8);
        \fill[red!55]     (3.2,0) rectangle (3.8,1.4);
        \fill[purple!55]  (4.2,0) rectangle (4.8,0.9);
        % Draw borders
        \draw[blue!80]   (0.2,0) rectangle (0.8,2.3);
        \draw[green!70]  (1.2,0) rectangle (1.8,0.6);
        \draw[orange!80] (2.2,0) rectangle (2.8,0.8);
        \draw[red!70]    (3.2,0) rectangle (3.8,1.4);
        \draw[purple!70] (4.2,0) rectangle (4.8,0.9);
        % Values
        \node[above,font=\tiny\bfseries] at (0.5,2.3)  {2.3с};
        \node[above,font=\tiny\bfseries] at (1.5,0.6)  {0.6с};
        \node[above,font=\tiny\bfseries] at (2.5,0.8)  {0.8с};
        \node[above,font=\tiny\bfseries] at (3.5,1.4)  {1.4с};
        \node[above,font=\tiny\bfseries] at (4.5,0.9)  {0.9с};
        % X labels
        \node[below,font=\tiny,align=center,text width=0.9cm] at (0.5,0)
          {Апп\\ };
        \node[below,font=\tiny,align=center,text width=0.9cm] at (1.5,0)
          {Fire-\\store};
        \node[below,font=\tiny,align=center,text width=0.9cm] at (2.5,0)
          {Node\\API};
        \node[below,font=\tiny,align=center,text width=0.9cm] at (3.5,0)
          {Claude};
        \node[below,font=\tiny,align=center,text width=0.9cm] at (4.5,0)
          {Groq};
      \end{tikzpicture}
    \end{column}
  \end{columns}
\end{frame}

%===============================================================================
\section{Шинэлэг тал}
%===============================================================================
\begin{frame}{Шинэлэг тал ба хувь нэмэр}
  \begin{columns}[T]
    \begin{column}{0.48\textwidth}
      \begin{block}{Техникийн шинэлэг тал}
        \begin{enumerate}
          \item \textbf{3 давхаргат архитектур:}\\Flutter — Node.js — Firebase/AWS
          \item \textbf{Үүлэн орчин:}\\Docker + AWS EC2/ECR
          \item \textbf{CI/CD:}\\GitHub Actions автоматжуулалт
          \item \textbf{AI чатбот:}\\Claude Haiku + Groq нөөцлөлт
          \item \textbf{Орон нутгийн хэрэгцээ:}\\fixed\_bookings, Dark Desert Theme
        \end{enumerate}
      \end{block}
    \end{column}
    \begin{column}{0.48\textwidth}
      \begin{alertblock}{Практик ач холбогдол}
        \begin{itemize}
          \item Даланзадгад хотын \textbf{5 спорт заал}-д нэвтрүүлэхэд бэлэн
          \item Иргэдийн цаг хугацааг хэмнэнэ
          \item Заалны ажилтнуудын ачааллыг бууруулна
          \item Захиалгын явц \textbf{ил тод} болно
        \end{itemize}
      \end{alertblock}
      \vspace{6pt}
      \begin{beamercolorbox}[sep=4pt,rounded=true]{block body}
        \small\centering
        Flutter + Firebase + Node.js + Docker + AWS + GitHub Actions + AI
        бүрэн стекийг Монголын орон нутагт хэрэглэсэн
        \textbf{анхны бүтээлүүдийн нэг}
      \end{beamercolorbox}
    \end{column}
  \end{columns}
\end{frame}

%===============================================================================
\section{Технологи ба консол холбоосууд}
%===============================================================================

%--- Tech overview table -------------------------------------------------------
\begin{frame}{Ашигласан бүх технологи — ерөнхий бүртгэл}
  \small
  \renewcommand{\arraystretch}{1.22}
  \begin{tabular}{llll}
    \toprule
    \rowcolor{navyblue}
    {\color{white}\textbf{Давхар}} &
    {\color{white}\textbf{Технологи}} &
    {\color{white}\textbf{Зориулалт}} &
    {\color{white}\textbf{Хувилбар}} \\
    \midrule
    \rowcolor{yellow!12}
    \multirow{5}{*}{\textbf{Flutter}} &
      flutter\_local\_notifications & Захиалгын сануулга (офлайн) & 18.0.0 \\
    \rowcolor{yellow!12}
      & firebase\_auth  & OTP утасны нэвтрэлт & 5.3.1 \\
    \rowcolor{yellow!12}
      & cloud\_firestore & Бодит цагийн өгөгдлийн сан & 5.4.4 \\
    \rowcolor{yellow!12}
      & cloud\_functions & Cloud Functions дуудах & 5.1.3 \\
    \rowcolor{yellow!12}
      & google\_fonts / http / intl & UI фонт, HTTP, огноо & latest \\
    \midrule
    \rowcolor{orange!10}
    \multirow{4}{*}{\textbf{Backend}} &
      express & REST API сервер & 4.18.2 \\
    \rowcolor{orange!10}
      & openai (Groq-д) & Groq Llama AI чат & 6.29.0 \\
    \rowcolor{orange!10}
      & firebase-admin & Firestore Admin SDK & 13.7.0 \\
    \rowcolor{orange!10}
      & @anthropic-ai/sdk & Claude Haiku API & 0.24.0 \\
    \midrule
    \rowcolor{blue!8}
    \multirow{3}{*}{\textbf{Cloud}} &
      Firebase Auth + Firestore & Нэвтрэлт, өгөгдлийн сан & v9 \\
    \rowcolor{blue!8}
      & Firebase Cloud Functions & Serverless Claude чат & v7.0.0 \\
    \rowcolor{blue!8}
      & AWS EC2 + ECR + Docker & Бэкэнд байршуулалт & ap-northeast-1 \\
    \midrule
    \rowcolor{green!8}
    \textbf{CI/CD} & GitHub Actions & Auto-deploy pipeline & v3 \\
    \bottomrule
  \end{tabular}
\end{frame}

%--- Firebase console ----------------------------------------------------------
\begin{frame}{Firebase консол — үр дүнг хаана харах вэ}
  \begin{columns}[T]
    \begin{column}{0.38\textwidth}
      \textbf{Console URL:}\\
      \small\texttt{console.firebase.google.com/\\project/\textbf{dz-zaal}}
      \vspace{6pt}

      \textbf{3 үйлчилгээ:}
      \begin{enumerate}
        \item \textbf{Authentication}\\→ Phone/OTP нэвтрэлт
        \item \textbf{Firestore Database}\\→ users, bookings,\\fixed\_bookings
        \item \textbf{Cloud Functions}\\→ chat function логууд
      \end{enumerate}
    \end{column}
    \begin{column}{0.60\textwidth}
      % Firebase console mockup
      \begin{tikzpicture}[scale=0.75]
        % Browser bar
        \fill[gray!20] (0,5.8) rectangle (10,6.4);
        \node[font=\tiny, gray!60, left] at (9.8,6.1)
          {\texttt{console.firebase.google.com/project/dz-zaal}};
        % Header
        \fill[black!85] (0,5.2) rectangle (10,5.8);
        \node[font=\small, white] at (2.5,5.5) {\textbf{Firebase} \quad dz-zaal};
        % Left sidebar
        \fill[gray!10] (0,0) rectangle (2.2,5.2);
        \foreach \y/\lbl in {4.7/Authentication, 4.2/Firestore Database,
                              3.7/Functions, 3.2/Hosting, 2.7/Storage} {
          \node[font=\tiny, black!70, right] at (0.1,\y) {\lbl};
          \draw[gray!30] (0,\y-0.25) -- (2.2,\y-0.25);
        }
        % Highlight Authentication
        \fill[blue!15] (0,4.45) rectangle (2.2,4.7);
        % Main content area
        \fill[white] (2.2,0) rectangle (10,5.2);
        % Table header
        \fill[gray!15] (2.3,4.7) rectangle (9.9,5.1);
        \node[font=\tiny, black!70] at (4.5,4.9) {Identifier};
        \node[font=\tiny, black!70] at (7.5,4.9) {Provider};
        \node[font=\tiny, black!70] at (9.2,4.9) {Created};
        % Rows (manual — foreach-д / тэмдэгт асуудал гардаг)
        \node[font=\tiny] at (4.5,4.3) {+976 9900 XXXX};
        \node[font=\tiny, green!60!black] at (7.2,4.3) {Phone};
        \node[font=\tiny, gray!60] at (9.2,4.3) {2026-03-20};
        \draw[gray!20] (2.2,4.1) -- (9.9,4.1);
        \node[font=\tiny] at (4.5,3.9) {+976 8811 XXXX};
        \node[font=\tiny, green!60!black] at (7.2,3.9) {Phone};
        \node[font=\tiny, gray!60] at (9.2,3.9) {2026-03-19};
        \draw[gray!20] (2.2,3.7) -- (9.9,3.7);
        \node[font=\tiny] at (4.5,3.5) {+976 9955 XXXX};
        \node[font=\tiny, green!60!black] at (7.2,3.5) {Phone};
        \node[font=\tiny, gray!60] at (9.2,3.5) {2026-03-18};
        \draw[gray!20] (2.2,3.3) -- (9.9,3.3);
        \node[font=\tiny] at (4.5,3.1) {+976 7700 XXXX};
        \node[font=\tiny, green!60!black] at (7.2,3.1) {Phone};
        \node[font=\tiny, gray!60] at (9.2,3.1) {2026-03-17};
        \draw[gray!20] (2.2,2.9) -- (9.9,2.9);
        \node[font=\tiny, navyblue] at (6.0,0.3)
          {\textbf{OTP-р нэвтэрсэн хэрэглэгчдийн жагсаалт}};
      \end{tikzpicture}
    \end{column}
  \end{columns}
\end{frame}

\begin{frame}{Firestore Database — захиалгын өгөгдөл}
  \begin{columns}[T]
    \begin{column}{0.36\textwidth}
      \textbf{Console URL:}\\
      \small\texttt{.../firestore/data}
      \vspace{8pt}

      \textbf{Харах боломжтой зүйлс:}
      \begin{itemize}
        \item \texttt{users} — бүртгэлтэй хэрэглэгч
        \item \texttt{bookings} — бүх захиалга\\
          status: \textcolor{okgreen}{active} /\\
          \textcolor{gray}{completed} /\\
          \textcolor{badred}{cancelled}
        \item \texttt{fixed\_bookings} —\\
          байгууллагын тогтмол цаг
      \end{itemize}
      \vspace{4pt}
      \begin{beamercolorbox}[sep=3pt,rounded=true]{block body}
        \tiny Захиалгын код:\\\textbf{DBK-} + 5 оронтой тоо
      \end{beamercolorbox}
    \end{column}
    \begin{column}{0.62\textwidth}
      % Firestore mockup
      \begin{tikzpicture}[scale=0.74]
        \fill[gray!20] (0,5.8) rectangle (11,6.3);
        \node[font=\tiny, gray!60] at (5.5,6.05)
          {\texttt{console.firebase.google.com/project/dz-zaal/firestore/data}};
        \fill[white] (0,0) rectangle (11,5.8);
        % 3 column panels
        \fill[gray!8]  (0,0)   rectangle (2.5,5.8);
        \fill[gray!5]  (2.6,0) rectangle (5.4,5.8);
        \fill[white]   (5.5,0) rectangle (11,5.8);
        % Column headers
        \fill[navyblue!80] (0,5.3) rectangle (2.5,5.8);
        \fill[navyblue!60] (2.6,5.3) rectangle (5.4,5.8);
        \fill[navyblue!40] (5.5,5.3) rectangle (11,5.8);
        \node[font=\tiny,white] at (1.25,5.55) {\textbf{Collections}};
        \node[font=\tiny,white] at (4.0,5.55)  {\textbf{Documents}};
        \node[font=\tiny,white] at (8.25,5.55) {\textbf{Fields}};
        % Collections
        \foreach \y/\coll/\active in {
          5.0/users/1, 4.6/bookings/0, 4.2/fixed\_bookings/0} {
          \ifnum\active=1
            \fill[blue!20] (0,\y-0.18) rectangle (2.5,\y+0.18);
          \fi
          \node[font=\tiny] at (1.25,\y) {\texttt{\coll}};
        }
        % Documents (bookings)
        \foreach \y/\doc in {
          5.0/DBK-48271, 4.6/DBK-33904, 4.2/DBK-71053,
          3.8/DBK-29447, 3.4/DBK-10823} {
          \node[font=\tiny] at (4.0,\y) {\texttt{\doc}};
          \draw[gray!25] (2.6,\y-0.2) -- (5.4,\y-0.2);
        }
        % Fields for selected document
        \node[font=\tiny, gray!60] at (8.25,4.8)
          {\texttt{bookingCode: "DBK-48271"}};
        \node[font=\tiny, gray!60] at (8.25,4.4)
          {\texttt{userId: "abc123uid"}};
        \node[font=\tiny, gray!60] at (8.25,4.0)
          {\texttt{venueId: "goviin\_arena"}};
        \node[font=\tiny, gray!60] at (8.25,3.6)
          {\texttt{date: "2026-03-24"}};
        \node[font=\tiny, gray!60] at (8.25,3.2)
          {\texttt{timeSlot: "10:00-11:00"}};
        \node[font=\tiny, okgreen] at (8.25,2.8)
          {\texttt{status: "active"}};
        \node[font=\tiny, gray!60] at (8.25,2.4)
          {\texttt{createdAt: Mar 24, 2026}};
      \end{tikzpicture}
    \end{column}
  \end{columns}
\end{frame}

%--- AWS console ---------------------------------------------------------------
\begin{frame}{AWS консол — EC2 ба ECR}
  \begin{columns}[T]
    \begin{column}{0.38\textwidth}
      \textbf{EC2 Console URL:}\\
      \small\texttt{ap-northeast-1.console.\\aws.amazon.com/ec2}
      \vspace{6pt}

      \textbf{EC2 Instance:}
      \begin{itemize}
        \item Instance ID: \texttt{i-0abc...}
        \item IP: \texttt{13.112.91.27}
        \item Region: ap-northeast-1 (Tokyo)
        \item Status: \textcolor{okgreen}{\textbf{running}}
        \item Port 3000 нээлттэй
      \end{itemize}
      \vspace{4pt}
      \textbf{ECR Repository:}
      \begin{itemize}
        \item Repo: \texttt{dz-zaal}
        \item Docker images SHA tag-тай
      \end{itemize}
    \end{column}
    \begin{column}{0.60\textwidth}
      \begin{tikzpicture}[scale=0.75]
        % AWS header
        \fill[black!90] (0,5.6) rectangle (10,6.2);
        \node[font=\small, white] at (1.2,5.9) {\textbf{aws}};
        \node[font=\tiny, white] at (4.0,5.9) {EC2 $>$ Instances};
        % EC2 table
        \fill[white] (0,0) rectangle (10,5.6);
        \fill[gray!12] (0,5.1) rectangle (10,5.5);
        \node[font=\tiny,black!70] at (1.0,5.3) {Name};
        \node[font=\tiny,black!70] at (3.5,5.3) {Instance ID};
        \node[font=\tiny,black!70] at (6.0,5.3) {State};
        \node[font=\tiny,black!70] at (8.0,5.3) {Public IP};
        % Running instance row
        \fill[green!6] (0,4.6) rectangle (10,5.0);
        \node[font=\tiny] at (1.0,4.8) {dz-zaal-backend};
        \node[font=\tiny] at (3.5,4.8) {\texttt{i-0a3f...d8c}};
        \fill[okgreen!20,rounded corners=2pt] (5.2,4.62) rectangle (6.8,4.98);
        \node[font=\tiny,okgreen] at (6.0,4.8)
          {$\bullet$ running};
        \node[font=\tiny] at (8.0,4.8) {\texttt{13.112.91.27}};
        \draw[gray!25] (0,4.6) -- (10,4.6);
        % ECR section below
        \fill[gray!8] (0,0) rectangle (10,3.5);
        \fill[black!90] (0,3.2) rectangle (10,3.5);
        \node[font=\tiny, white] at (3.0,3.35)
          {ECR $>$ Repositories $>$ dz-zaal};
        \fill[gray!12] (0,2.8) rectangle (10,3.1);
        \node[font=\tiny,black!70] at (1.5,2.95) {Image tag};
        \node[font=\tiny,black!70] at (5.0,2.95) {Pushed at};
        \node[font=\tiny,black!70] at (8.0,2.95) {Size};
        \node[font=\tiny] at (1.5,2.5) {\texttt{a3f7d8c2}};
        \node[font=\tiny,gray!60] at (5.0,2.5) {2026-03-24 14:22};
        \node[font=\tiny,gray!60] at (8.0,2.5) {48.3 MB};
        \draw[gray!20] (0,2.3) -- (10,2.3);
        \node[font=\tiny] at (1.5,2.1) {\texttt{b1e9f3a4}};
        \node[font=\tiny,gray!60] at (5.0,2.1) {2026-03-22 09:15};
        \node[font=\tiny,gray!60] at (8.0,2.1) {47.9 MB};
        \draw[gray!20] (0,1.9) -- (10,1.9);
        \node[font=\tiny] at (1.5,1.7) {\texttt{c4d2b8e1}};
        \node[font=\tiny,gray!60] at (5.0,1.7) {2026-03-20 16:40};
        \node[font=\tiny,gray!60] at (8.0,1.7) {47.8 MB};
        \draw[gray!20] (0,1.5) -- (10,1.5);
      \end{tikzpicture}
    \end{column}
  \end{columns}
\end{frame}

%--- GitHub Actions ------------------------------------------------------------
\begin{frame}{GitHub Actions — CI/CD workflow}
  \begin{columns}[T]
    \begin{column}{0.38\textwidth}
      \textbf{URL:}\\
      \small\texttt{github.com/[repo]/actions}
      \vspace{6pt}

      \textbf{Workflow файл:}\\
      \texttt{.github/workflows/\\deploy.yml}
      \vspace{6pt}

      \textbf{Trigger:}
      \begin{itemize}
        \item \texttt{backend/} хавтаст\\push хийхэд
        \item \texttt{main} branch дээр
      \end{itemize}
      \vspace{4pt}
      \textbf{Secrets хадгалах:}
      \begin{itemize}
        \item \texttt{AWS\_ACCESS\_KEY\_ID}
        \item \texttt{AWS\_SECRET\_ACCESS\_KEY}
        \item \texttt{EC2\_SSH\_KEY}
        \item \texttt{EC2\_HOST}
      \end{itemize}
    \end{column}
    \begin{column}{0.60\textwidth}
      \begin{tikzpicture}[scale=0.75]
        % GitHub header
        \fill[black!88] (0,5.8) rectangle (10,6.4);
        \node[font=\small, white] at (1.5,6.1)
          {\textbf{GitHub Actions}};
        \node[font=\tiny, white] at (6.0,6.1)
          {dz-zaal / .github / workflows / deploy.yml};
        % Content
        \fill[white] (0,0) rectangle (10,5.8);
        \fill[gray!10] (0,5.3) rectangle (10,5.8);
        \node[font=\small, black!80] at (2.0,5.55)
          {\textbf{Deploy Backend}};
        % Workflow runs (manual rows)
        \fill[okgreen!10] (0.1,4.62) rectangle (9.9,5.18);
        \fill[okgreen!70,rounded corners=2pt] (0.15,4.72) rectangle (0.65,5.08);
        \node[font=\small,white] at (0.4,4.9) {$\checkmark$};
        \node[font=\tiny,right] at (0.7,5.02) {backend: AI endpoint fix};
        \node[font=\tiny,gray!60,right] at (0.7,4.74) {\texttt{a3f7d8c} \enspace Mar 24 14:20 \enspace 3 min};
        \draw[gray!20] (0,4.62) -- (10,4.62);

        \fill[okgreen!10] (0.1,4.02) rectangle (9.9,4.58);
        \fill[okgreen!70,rounded corners=2pt] (0.15,4.12) rectangle (0.65,4.48);
        \node[font=\small,white] at (0.4,4.3) {$\checkmark$};
        \node[font=\tiny,right] at (0.7,4.42) {backend: add cron cleanup};
        \node[font=\tiny,gray!60,right] at (0.7,4.14) {\texttt{b1e9f3a} \enspace Mar 22 09:12 \enspace 3 min};
        \draw[gray!20] (0,4.02) -- (10,4.02);

        \fill[okgreen!10] (0.1,3.42) rectangle (9.9,3.98);
        \fill[okgreen!70,rounded corners=2pt] (0.15,3.52) rectangle (0.65,3.88);
        \node[font=\small,white] at (0.4,3.7) {$\checkmark$};
        \node[font=\tiny,right] at (0.7,3.82) {backend: update groq model};
        \node[font=\tiny,gray!60,right] at (0.7,3.54) {\texttt{c4d2b8e} \enspace Mar 20 16:38 \enspace 4 min};
        \draw[gray!20] (0,3.42) -- (10,3.42);

        \fill[okgreen!10] (0.1,2.82) rectangle (9.9,3.38);
        \fill[okgreen!70,rounded corners=2pt] (0.15,2.92) rectangle (0.65,3.28);
        \node[font=\small,white] at (0.4,3.1) {$\checkmark$};
        \node[font=\tiny,right] at (0.7,3.22) {backend: fix CORS headers};
        \node[font=\tiny,gray!60,right] at (0.7,2.94) {\texttt{d8e1c3f} \enspace Mar 18 11:05 \enspace 3 min};
        \draw[gray!20] (0,2.82) -- (10,2.82);

        \fill[okgreen!10] (0.1,2.22) rectangle (9.9,2.78);
        \fill[okgreen!70,rounded corners=2pt] (0.15,2.32) rectangle (0.65,2.68);
        \node[font=\small,white] at (0.4,2.5) {$\checkmark$};
        \node[font=\tiny,right] at (0.7,2.62) {backend: initial deploy};
        \node[font=\tiny,gray!60,right] at (0.7,2.34) {\texttt{e2f4a1b} \enspace Mar 15 09:00 \enspace 4 min};
        \draw[gray!20] (0,2.22) -- (10,2.22);

        % Step detail at bottom
        \fill[gray!6] (0,0) rectangle (10,1.9);
        \node[font=\tiny,black!70] at (5,1.72)
          {\textbf{Сүүлийн run-ийн алхмууд:}};
        % Steps manually
        \fill[okgreen!20,rounded corners=2pt] (0.1,0.2) rectangle (1.8,1.3);
        \node[font=\tiny,okgreen] at (0.95,1.05) {$\checkmark$};
        \node[font=\tiny,align=center] at (0.95,0.62) {Checkout};

        \fill[okgreen!20,rounded corners=2pt] (2.0,0.2) rectangle (4.0,1.3);
        \node[font=\tiny,okgreen] at (3.0,1.05) {$\checkmark$};
        \node[font=\tiny,align=center] at (3.0,0.62) {AWS\\Credentials};

        \fill[okgreen!20,rounded corners=2pt] (4.2,0.2) rectangle (6.0,1.3);
        \node[font=\tiny,okgreen] at (5.1,1.05) {$\checkmark$};
        \node[font=\tiny,align=center] at (5.1,0.62) {ECR Login};

        \fill[okgreen!20,rounded corners=2pt] (6.2,0.2) rectangle (8.4,1.3);
        \node[font=\tiny,okgreen] at (7.3,1.05) {$\checkmark$};
        \node[font=\tiny,align=center] at (7.3,0.62) {Docker\\Build+Push};

        \fill[okgreen!20,rounded corners=2pt] (8.6,0.2) rectangle (10.0,1.3);
        \node[font=\tiny,okgreen] at (9.3,1.05) {$\checkmark$};
        \node[font=\tiny,align=center] at (9.3,0.62) {SSH\\Deploy};
      \end{tikzpicture}
    \end{column}
  \end{columns}
\end{frame}

%--- AI consoles ---------------------------------------------------------------
\begin{frame}{AI API консолууд — Groq ба Anthropic}
  \begin{columns}[T]
    \begin{column}{0.48\textwidth}
      \begin{block}{Groq Console}
        \textbf{URL:} \texttt{console.groq.com}\\[4pt]
        \textbf{Model:} \texttt{llama-3.1-8b-instant}\\
        \textbf{API base:} \texttt{api.groq.com/openai/v1}\\
        \textbf{SDK:} OpenAI-compatible (\texttt{openai} ^6.29)\\[4pt]
        \textbf{Консолд харагдах:}
        \begin{itemize}
          \item API key удирдах
          \item Request/token usage лог
          \item Rate limits (6000 req/min)
          \item Latency: \textbf{дундаж 0.9с}
        \end{itemize}
        \vspace{4pt}
        \begin{beamercolorbox}[sep=3pt,rounded=true]{block body}
          \tiny\textbf{Backend server.js:}\\
          \texttt{baseURL: 'https://api.groq.com/openai/v1'}\\
          \texttt{apiKey: process.env.GROQ\_API\_KEY}
        \end{beamercolorbox}
      \end{block}
    \end{column}
    \begin{column}{0.48\textwidth}
      \begin{block}{Anthropic Console}
        \textbf{URL:} \texttt{console.anthropic.com}\\[4pt]
        \textbf{Model:} \texttt{claude-haiku-4-5-20251001}\\
        \textbf{SDK:} \texttt{@anthropic-ai/sdk} ^0.78\\
        \textbf{Deploy:} Firebase Cloud Functions\\[4pt]
        \textbf{Консолд харагдах:}
        \begin{itemize}
          \item API key удирдах
          \item Token usage (input/output)
          \item Latency: \textbf{дундаж 1.4с}
          \item Cost per 1M tokens
        \end{itemize}
        \vspace{4pt}
        \begin{beamercolorbox}[sep=3pt,rounded=true]{block body}
          \tiny\textbf{functions/index.js:}\\
          \texttt{model: 'claude-haiku-4-5-20251001'}\\
          \texttt{secret: CLAUDE\_API\_KEY}
        \end{beamercolorbox}
      \end{block}
    \end{column}
  \end{columns}
\end{frame}

%--- Flutter packages ----------------------------------------------------------
\begin{frame}{Flutter / Dart package-ууд — pub.dev}
  \begin{center}
  \small
  \renewcommand{\arraystretch}{1.2}
  \begin{tabular}{llll}
    \toprule
    \rowcolor{navyblue}
    {\color{white}\textbf{Package}} &
    {\color{white}\textbf{Хувилбар}} &
    {\color{white}\textbf{Зориулалт}} &
    {\color{white}\textbf{pub.dev URL}} \\
    \midrule
    \rowcolor{rowgray}
    firebase\_auth         & 5.3.1  & OTP утасны нэвтрэлт   & \texttt{.../firebase\_auth} \\
    cloud\_firestore       & 5.4.4  & Өгөгдлийн сан (realtime) & \texttt{.../cloud\_firestore} \\
    \rowcolor{rowgray}
    cloud\_functions       & 5.1.3  & Cloud Functions дуудах & \texttt{.../cloud\_functions} \\
    firebase\_core         & 3.6.0  & Firebase SDK суурь    & \texttt{.../firebase\_core} \\
    \rowcolor{rowgray}
    flutter\_local\_notifications & 18.0.0 & Офлайн сануулга & \texttt{.../flutter\_local...} \\
    timezone               & 0.9.4  & Цагийн бүс (сануулганд) & \texttt{.../timezone} \\
    \rowcolor{rowgray}
    google\_fonts          & 6.1.0  & Montserrat фонт       & \texttt{.../google\_fonts} \\
    http                   & 1.2.0  & REST API дуудах       & \texttt{.../http} \\
    \rowcolor{rowgray}
    intl                   & 0.19.0 & Огноо форматлах       & \texttt{.../intl} \\
    cupertino\_icons       & 1.0.8  & iOS дүрсүүд           & \texttt{.../cupertino\_icons} \\
    \bottomrule
  \end{tabular}
  \end{center}
  \vspace{2pt}
  \begin{beamercolorbox}[sep=3pt,center,rounded=true]{block body}
    \small Бүх package: \texttt{https://pub.dev/packages/[package-name]}
  \end{beamercolorbox}
\end{frame}

%===============================================================================
\section{Дүгнэлт}
%===============================================================================
\begin{frame}{Дүгнэлт}
  \begin{columns}[T]
    \begin{column}{0.52\textwidth}
      \begin{block}{Хэрэгжүүлсэн бүтээгдэхүүн}
        \begin{itemize}
          \item[\YES] Flutter мобайл апп (Android + iOS)
          \item[\YES] Firebase Auth (OTP) + Cloud Firestore
          \item[\YES] Node.js/Express REST API бэкэнд
          \item[\YES] Docker + AWS EC2 байршуулалт
          \item[\YES] GitHub Actions CI/CD хоолой
          \item[\YES] Claude Haiku + Groq AI чатбот
          \item[\YES] 10/10 функциональ тест тэнцсэн
        \end{itemize}
      \end{block}
      \vspace{4pt}
      \textbf{Бүх 7 зорилт биелэгдсэн.}
    \end{column}
    \begin{column}{0.45\textwidth}
      \begin{block}{Цаашдын хөгжлийн чиглэл}
        \begin{itemize}
          \item QPay / SocialPay төлбөрийн систем
          \item Push мэдэгдэл (FCM)
          \item Захирагчийн admin панел
          \item GPS байршлын функц
          \item Статистик, тайлан үүсгэх
          \item iOS App Store, Google Play-д нийтлэх
        \end{itemize}
      \end{block}
    \end{column}
  \end{columns}
\end{frame}

%--- Thank you -----------------------------------------------------------------
\begin{frame}[plain]
  \vspace{-0.3cm}
  {\color{navyblue}\hrule height 2.5pt}
  \vspace{0.5cm}
  \begin{center}
    {\Huge\bfseries\color{navyblue} Анхааралтай сонссонд}\\[6pt]
    {\Huge\bfseries\color{goldcolor} баярлалаа!}
  \end{center}
  \vspace{0.4cm}
  {\color{goldcolor}\hrule height 1.5pt}
  \vspace{0.4cm}
  \begin{center}
    \begin{beamercolorbox}[sep=0.6em,center,rounded=true]{block title}
      \large Асуулт хүлээн авах боломжтой
    \end{beamercolorbox}
  \end{center}
  \vspace{0.3cm}
  \begin{center}
    \small
    \begin{tabular}{rl}
      \textbf{Оюутан:} & М.Түвшинжаргал (В221940619) \\
      \textbf{Тэнхим:} & Мэдээллийн технологийн тэнхим, ШУТИС-МХТС \\
      \textbf{Жил:}    & 2026 он
    \end{tabular}
  \end{center}
  \vspace{0.3cm}
  {\color{navyblue}\hrule height 2.5pt}
\end{frame}

\end{document}
